% Generated draft: requires mathematical and source review.
\documentclass[11pt]{article}
\usepackage[margin=25mm]{geometry}
\usepackage{fontspec}
\setmainfont{Latin Modern Roman}
\usepackage{amsmath,amssymb,mathtools,unicode-math}
\setmathfont{Latin Modern Math}
\usepackage{enumitem,needspace,xurl,hyperref}
\setlength{\parindent}{0pt}
\setlength{\parskip}{5pt}
\setlength{\emergencystretch}{3em}
\newcommand{\N}{\mathbb{N}}
\newcommand{\Z}{\mathbb{Z}}
\newcommand{\Q}{\mathbb{Q}}
\newcommand{\R}{\mathbb{R}}
\newcommand{\C}{\mathbb{C}}
\newcommand{\F}{\mathbb{F}}
\newcommand{\A}{\mathbb{A}}
\newcommand{\PP}{\mathbb{P}}
\newcommand{\E}{\mathbb{E}}
\newcommand{\eps}{\varepsilon}
\newcommand{\dd}{\,\mathrm{d}}
\begin{document}
\section*{UnsolvedMath ID: 30001462}
\textbf{Source problem code:} catalogue entry
\subsection{Definitions and mathematical statement}
The supplied record refers to an object \(D_h\), its defining fan, and quantities \(d_i\), but does not define these objects, the parameter \(h\), the index range for \(i\), or the ambient space. Consequently, the following is a conditional transcription of the intended question, rather than a fully specified mathematical conjecture.

Write \(\Sigma_h\) for the fan referred to in the source and \(I\) for the intended index set of the quantities \(d_i\). For cones \(F\) and \(\tau\), write \(F\prec_1\tau\) when \(F\) is a face of \(\tau\) with \(\dim F=\dim\tau-1\); such a face is called a facet. Interpret two cones meeting in a common facet to mean that their intersection is a facet of each. These conventions explicate the terminology; the missing source context must confirm their applicability.

\textbf{Question as supplied.} For every cone \(\sigma\in\Sigma_h\), do there exist a positive integer \(r\), cones \(\sigma_1,\ldots,\sigma_r\in\Sigma_h\), an index \(i\in I\), and a facet \(F\prec_1\sigma_r\) such that
\[
\sigma_1=\sigma,
\qquad
\sigma_j\cap\sigma_{j+1}\prec_1\sigma_j
\quad\text{and}\quad
\sigma_j\cap\sigma_{j+1}\prec_1\sigma_{j+1}
\quad(1\leq j<r),
\]
and \(F\) is defined by \(d_i=0\)? Under the usual interpretation of the last phrase as a defining hyperplane section, its intended expression is
\[
F=\sigma_r\cap\{x:d_i(x)=0\},
\]
where the meaning and domain of \(d_i(x)\) must be supplied by the original context. For \(r=1\), there are no consecutive-cone conditions.

\textbf{Unresolved qualifications.} The record says every cone, without restricting to maximal cones or cones of a specified dimension. Under the facet convention above, consecutive cones necessarily have equal dimension. Moreover, a fan ordinarily contains the zero cone, which has no facet under this convention. The intended range of cones and any exceptional conventions therefore require clarification; no restriction to maximal cones has been silently imposed. The permitted values and quantification of \(h\), the construction of \(D_h\) and \(\Sigma_h\), and the meaning of the equations \(d_i=0\) remain unspecified.
\subsection{Short English statement}
The question asks whether every cone in the fan defining \(D_h\) can reach a cone with a facet defined by some \(d_i=0\), through a finite sequence in which consecutive cones share a facet. The supplied record omits the definitions and parameter ranges needed for a complete statement, and its quantifier over every cone requires clarification. This is a statement-only draft requiring human mathematical review.
\subsection{Sources}
\begin{itemize}
\item Local UnsolvedMath registry: permanent ID \(30001462\), source problem code \texttt{catalogue entry}, title \emph{Facet Connectivity in the Fan of \(D_h\)}; collection: Oberwolfach Reports Open Problems. The supplied record gives the proposed year as 2010 and identifies no individual proposer.
\item Supplied report citation: \emph{Algebraic Groups} (2010), Oberwolfach Reports, pp.\ 1101--1163, DOI \texttt{10.4171/owr/2010/19}. Source URL and supplied literature-source link, labeled ``original OWR report'': \url{https://doi.org/10.4171/owr/2010/19}. The supplied background describes Oberwolfach Reports as documenting mathematical workshops and their open-problem discussions. No precise page or contribution containing this question is identified.
\item Supplied statement-verification metadata labels the statement \texttt{corrected\_verified} and reports that its mathematical content was checked against the official report, with notation, context, and question form normalized without changing the source-level claim. This is a supplied verification claim, not an independent verification performed for this draft.
\item Supplied literature assessment, checked on \texttt{2026-08-21}: status labeled open. The assessment says that the source labels facet-connectivity of this fan as the main open problem and that no later primary result located by the cutoff establishes the requested connectivity in the stated generality. Its listed checked source is the original report linked above. This assessment is supplied dated curation and literature triage, not independently verified; time-sensitive status and cited claims require independent checking.
\end{itemize}
\end{document}
